\documentclass[12pt]{article}
\usepackage{color,soul}
% Import preambles and macros for notes
\input{../../../../preambles/notes-preamble.tex}
\input{../../../../preambles/notes-macros.tex}
\input{../../../../preambles/theorem-system-section.tex}

\title{MATH 420 Lecture 15}
\author{Deepak Jassal}
\date{March 9\textsuperscript{th}, 2026}

\begin{document}
\maketitle
\stepcounter{section}
\defn{}{
    A ring homomorphism $i_B:A\to B$ is said to be \underline{finite-type}, and $B$ a \underline{finitely generated} $A$-algebra, if $B$ is generated by a finite set of elements as an $A$-algebra. 
}
Compare this definition to the definition of finite-dimensional vector space $V$ over a field $F$. That being
\[
    V=\spn_F\{v_1,\dots,v_k\}
\]
for some $v_1,\dots,v_k\in V$.\\
Interpreting this for an $A$-algebra, we want:\\
there exists $b_1,\dots,b_k\in B$ such that for all $b\in B$, there exists a polynomial $P_b\in A[x_1,\dots,x_k]$ such that $b=i_B(P_b)(b_1,\dots,b_k)$.\\
Unlike vector spaces, the size of a minimal generating set (i.e., a generating fet $S$ such that $S\setminus\{s\}$ is not a generating set for all $s\in S$) of a general $A$-algebra $B$ is \underline{not} the same. So ``dimension'' is not well defined.\\
A natural family of algebras (they can be thought of as general $A$-algebras, of $F$-algebras, for $F=\Z,\Q,\R,\C$,etc.) are so-called \underline{rings of polynomial} invariants/covariants. If we want to study polynomials, we can consider how htey behave when acted upon by a group.\\
Consdier the group
\[
    GL_2(R)=\left\{\begin{bmatrix}
        a_{11}&a_{12}\\a_{21}&a_{22}
    \end{bmatrix}:
    a_{11}a_{22}-a_{12}a_{21}\neq0\right\}
\]
with $R$ a unital, commutative ring. Let 
\[
    P_2(R)=\{f_2x^2+f_1xy+f_0y^2:f_2,f_1,f_0\in R\}.
\]
$GL_2(R)$ acts on $P_2(R)$ via the \underline{substitution action}:
\[
    g=
    \begin{bmatrix}
        a_{11}&a_{12}\\a_{21}&a_{22}
    \end{bmatrix}\to f\in P_2(R),
\]
by 
\[
    g\cdot f=f(a_{11}x+a_{21}y,a_{12}x+a_{22}y)=f_2(a_{11}x+a_{21}y)^2+f_1(a_{11}x+a_{21}y)(a_{12}x+a_{22}y)+f_0(a_{12}x+a_{22}y)^2.
\]
Restrict to the subgroup $G=GL_2^{\pm1}(R)$ ($\det=\pm1$).\\
Q: does there exists a \underline{non-zero} polynomial, $Q(f)=Q(f_2,f_1,f_0)$ such that for all $g\in G$, $Q(g\cdot f)=Q(f)$?\\
Yes, $Q$ can be taken to be the \underline{discriminant}, $\Delta(f)=f_1^2-4f_2f_0$.\\
Let's try this for $G=GL_2(\Z)=GL_2^{\pm1}(\Z)$. 
\[
    G=\left\langle
    \begin{bmatrix}
        1&1\\0&1
    \end{bmatrix},
    \begin{bmatrix}
        0&1\\-1&0
    \end{bmatrix},
    \begin{bmatrix}
        0&1\\1&0
    \end{bmatrix}
    \right\rangle=\langle E_1,S,T\rangle.
\] 
\[
    E_1\cdot f=f_2x^2+f_1(x)(x+y)+f_0(x+y)^2=(f_2+f_1+f_0)x^2+(f_1+2f_0)xy+f_0y^2,
\]
\[
    \Delta(E_1\cdot f)=(f_1+2f_0)^2-4(f_2+f_1+f_0)f_0=f_1^2-4f_2f_0=\Delta(f).
\]
\[
    S\cdot f=f_2(-y^2)+f_1(-y)(x)+f_0x^2=f_0x^2-f_1xy+f_2y^2
\]
\[
    \Delta(S\cdot f)=(-f_1)^2-4f_0f_2=f_1^2-4f_2f_0=\Delta(f).
\]
Inv$_2(\Z)$ si the set of polynomials $Q(F)$ such that $Q(g\cdot f)=Q(f)$ for all $g\in G$. We just showed that $\langle\Delta(f)\rangle\subseteq$Inv$_2(\Z)$ (they are actually equal).
\[
    P_3(R)=\{f_3x^3+f_2x^2y+f+1xy^2+f_0y^3\}.
\]
$G$ acts on $P_3(R)$ by substitution. 
\[
    \Delta_3(f)=f_2^2f_1^2-4f_3f_1^3-4f_0f_2^3-27f_3^2f_0^2+18f_3f_2f_1f_0.
\]
$\Delta_3(g\cdot f)=\Delta_3(f)$ for all $g\in G$. Inv$_3(\Z)=\langle\Delta_3(f)\rangle$.\\
However for quartics we have, Inv$_4(\Z)=\langle I(f),J(f)\rangle$, with $I(f)=12f_4f_0-3f_3f_1+f_2^2$ and $J(f)=72f_4f_2f_0+9f_3f_2f_1-27f_4f_1^2-27f_0f_3^2-2f_2^3$.\\
No longer true for quintics. 
\[
    \text{Inv}_5=\langle H_6,h_8,H_{10},H_{15}\rangle/ n
\]
\[
    H^2_{15}=P(H_6,H_8,H_{10})
\]
\end{document}